\cotefile{14}
\scanpage{45}
\mspage{2.}
$i\handwritten{\geqslant} n$, surjectif pour $i=n-1$. En fait, sauf erreur je sais prouver la conj.$A$ à partir de la forme plus faible suivante (dans la nature d'un théorème de Lefschetz ``faible'') :

\uline{Conjecture} $A'$. Pour tout entier $i$ tel que $2i\handwritten{\geqslant} n-1$, l'application image directe $C^i(Y)\handwritten{\longrightarrow} C^{i+1}(X)$ est surjective, i.e. tout cycle algébrique de dimension $j\handwritten{<}n/2$ sur $X$ est $\handwritten{\tau}$-équivalent à l'image d'un cycle algébrique à coeff. rationnels de $Y$.

Ici encore il suffit sauf erreur de faire $n=2i+1$. D'ailleurs, dans l'argument utilisé pour ramener $A$ à $A'$, il semble qu'il faille démontrer $A'$ sur un corps de base pas nécessairement algébriquement clos, et pour des classes de cycles rationnelles sur ce dernier, \uline{ou} se borner à énoncer $A$ pour l'équivalence cohomologique (mais alors $A$ perd le caractère ``purement géométrique'' que je t'avais promis !).

Comme je te l'avais dit, $A$ implique ``la formule de Künneth pour les cycles'' (j'avais d'ailleurs charrié en prétendant que les deux énoncés sont équivalents). Par là, il implique tous les théorèmes d'entiers qu'on peut désirer (par exemple pour des coefficients de polynômes caractéristiques), à cela près qu'il semble possible qu'il y ait des puissances de $p=\operatorname{car} k$ en dénominateur. Cela implique donc les conjectures de Weil, à la seule exception de la question des valeurs absolues des valeurs propres, qui sera justiciable de $B$. Je te signale aussi que c'est la conjecture $A$ qui semble le minimum minimorum pour donner une définition en forme utilisable de la notion de motif sur un corps.

\uline{Conjecture} $B$\quad Supposons $n=2m$, et soit $P^m(X)$ le noyau de l'homomorphisme de multiplication par $\handwritten{\xi}$ de $C^m(X)$ dans $C^{m+1}(X)$. Alors la forme $(-1)^m\handwritten{\unsure{\varepsilon|\mathcal{E}}}(xx')$ sur $P^m(X)$ est définie positive.

Pour les conjectures de Weil, il suffirait d'ailleurs de prouver que cette forme est positive. Mais la forme forte implique aussi d'autres résultats agréables, comme le fait que $\handwritten{\tau}$-équivalence $=$ équivalence numérique ($=$ équivalence homologique à coeff $\underline{\mathrm{Q}}_{\unsure{l|1|\ell}}$, cette dernière étant coincée entre les deux) et le fait que les $C^m(X)$ sont de dimension finie,
% archive sheet number: 44
\scanpage{46}
\mspage{3.}
en fait, que les groupes de classes de cycles pour la $\handwritten{\tau}$-équivalence sont de type fini. (De là, on déduit d'ailleurs formellement que $C^m(X)\otimes_{\underline{\mathrm{Q}}}\underline{\mathrm{Q}}_{\unsure{l|1|\ell}}\handwritten{\longrightarrow} H^{2m}(X)(m)$ est injectif, donc que le rang de $C^m(X)$ est majoré par $b_{2m}(X)$). D'autre part, $A$ et $B$ impliquent également que la catégorie des motifs construits par des variétés projectives non singulières est semi-simple, et de façon plus terre à terre, que l'anneau des classes de correspondance algébriques (mod $\handwritten{\tau}$-équivalence et tensorisé par $\underline{\mathrm{Q}}$) de $X$ est une algèbre finie semi-simple sur $\underline{\mathrm{Q}}$, qu'on peut d'ailleurs munir d'une involution et d'une trace satisfaisant les conditions habituelles. On en conclut aussi le résultat suivant, qui peut être considéré dans une certaine mesure comme une généralisation de $A$ (lorsqu'on admet que la variante cohomologique de $A$ est déjà acquise) : Soit $H^{2i}(X)\handwritten{\longrightarrow} H^{2i+2r}(X')$ un homomorphisme défini par une classe de correspondance algébrique, et $C^i(X)\handwritten{\longrightarrow} C^{i+r}(X')$ l'homomorphisme qu'elle définit sur les cycles. Alors un élément de $C^{i+r}(X')$ est dans l'image de $C^i(X)$ si et seulement si son image dans $H^{2i+2r}$ est dans celle de $H^{2i}$. - Notes aussi que le résultat de semi-simplicité signalé plus haut impliquerait le résultat de semi-simplicité analogue pour les opérations du groupe de Galois, dans les conjectures de Tate. Enfin, je me suis à peu près convaincu que $A$ et $B$ impliquent également une théorie raisonnable des variétés abéliennes qui interviennent comme variétés de paramètres pour les familles continues de cycles algébriques, et notamment permettent d'obtenir les relations qu'il faut entre ces ``jacobiennes intermédiaires'' (qui peuvent être conçues comme les ``morceaux de nature algébrique'' des jacobiennes de Weil) et la cohomologie de dimension impaire. Mais il est fort possible que la démonstration de $B$ soit déjà liée à l'introduction de ces variétés abéliennes.

En tous cas, $A$ et $A'$ me semblent, dans un sens, préliminaires à $B$, et il semble raisonnable de commencer par elles. Je te signale la forme équivalente assez suggestive de $A'$ :

\uline{Conjecture} $A''$. Soit $U$ l'ouvert affine (de dimension $n$) $X-Y$, alors $C^i(U)=0$ pour tout $i$ tel que $2i\handwritten{>}n$.
% archive sheet number: 45
\scanpage{47}
\mspage{4.}
En fait, il est possible que ceci soit vrai pour toute variété affine lisse, et même que pour toute variété (affine ou non) lisse, tout cycle cohomologiquement équivalent à $0$ soit $\handwritten{\tau}$-équivalent à $0$ (pour les variétés projectives, cela résulterait, comme j'ai dit, de $A$ et $B$).

Ce qu'il faut pour le moment, c'est inventer un procédé pour déformer un cycle de dimension pas trop grande, pour le pousser à l'infini. Peut-être aurais-tu envie d'y réfléchir de ton côté ? Je viens seulement de m'y mettre, aujourd'hui même, et t'écris faute de trouver une idée.

Bien à toi\\
\hfill A. GROTHENDIECK\\

N.B. $H^i(X)$ désigne la cohomologie à coefficients $\mathbb{Q}_{\unsure{l|1|\ell}}$, éventuellement tordus par une puissance tensorielle de $T_l(\mathbb{G}\mathrm{m})\otimes_{\mathbb{Z}}\mathbb{Q}_l$, qui est omise dans les notations pour simplifier.
% archive sheet number: 46
\scanpage{48}
Sorites sur VA,\\
Picard, Albanese..... 
% archive sheet number: 47
\scanpage{49}
\uline{Résumé} On peut \unsure{exprimer|définir} des V.A. sur $k$ (\unsure{alg|alg.} clos fixé), \deleted{\unsure{des}} en termes de cycles algébriques \inserted{alg. équiv. à $0$} sur une variété, soit via les diviseurs ($\deleted{\unsure{A}}B_X$\inserted{$={}$ Picard \unsure{connexe|réduit}}) soit via les $0$-cycles ($A_X\deleted{\unsure{)}}={}$ Albanese). On obtient \deleted{donc}\inserted{alors}, pour deux variétés abéliennes $A$ et $B$, définies \inserted{chacune,} en termes de $X$ \deleted{resp.} de $Y$, de l'une ou de l'autre façon, et un cycle alg. de dim convenable $Z$ sur $X\times Y$, par $\unsure{\mathscr{D}|\delta|D}\unsure{\leadsto|-|\mapsto} Z(\unsure{\delta|\mathscr{D}|D})$, une application $A(k)\longrightarrow B(k)$, \deleted{\unsure{et une|et même}} (il y a env. $4=2\times2$ cas à distinguer). Nous montrons qu'elle est tjrs associée à un hom de V.A. $u_Z:A\longrightarrow B$, \inserted{\unsure{ne|en} dépendant que de la classe d'\unsure{équiv|équivalence} \uline{numérique} de $Z$.} \deleted{et \unsure{pour|que}}\inserted{D'ailleurs,} si on admet des cycles à coeff. rationnels, \deleted{\unsure{les|tous}} tout hom de V.A. $A\unsure{-|\longrightarrow} B$ s'obtient de cette façon. En termes de cohomologie $l$-adique, on peut dire que les hom $H^1(B)-H^1(A)$ associés à des $\mathrm{Q}$-hom de V.A. $A-B$, sont aussi ceux associés à des $\mathrm{Q}$-cycles
% archive sheet number: 48
\scanpage{50}
algébriques sur $X\times Y$, via les identifications \unsure{du|des} type (pour une variété $T$)
\[
\left\{\begin{aligned}
H^1(A_T,\mathrm{Q}_{\ell})&\xrightarrow{\sim}H^1(T,\mathrm{Q}_{\ell})\\
H^{\deleted{\unsure{2n-|2n-1}}1}(B_T,\mathrm{Q}_{\ell})&\xleftarrow{\sim}H^{2n-1}(T,\mathrm{Q}_{\ell}\deleted{\unsure{(}})(n-1)
\end{aligned}\right.
\]
De plus, l'hom correspondant entre groupes de coh. $l$-adiques relatifs à $X$ et à $Y$, est un isom ssi le $\mathrm{Q}$-hom \deleted{correspondant} envisagé $A\unsure{-|\longrightarrow}B$ est une \uline{isogénie}. Cela prouve en particulier que dans ce cas, il y a une classe de corr. $\mathrm{Q}$-algébriques sur $X\times Y$ définissant l'hom inverse sur la cohomologie.

Nous verrons que cette circonstance se produit en particulier si on a des classes de diviseurs amples $\unsure{\zeta|\xi}_1,\ldots,\zeta_{n-1}$ sur $X$ de dim $n$, en prenant l'hom de cup-produit par \deleted{$\unsure{\xi|\zeta}$} $\zeta_1\cdot\ldots\zeta_{n-1}:H^1(X,\mathrm{Q}_{\ell})-H^{2n-1}(X,\mathrm{Q}_{\ell})(n-1)$.
% archive sheet number: 49
\scanpage{51}
\margin{Cas transposés l'un de l'autre}
\diagram{\scriptsize
\begin{tabular}{@{}c|c@{}}
\hfill I & \hfill II\\
\begin{tikzcd}[ampersand replacement=\&,column sep=small,row sep=small]
H^1(Y)\arrow[r] \& H^1(X)\\
H^1(A_Y)\arrow[u,"\varphi_Y"']\arrow[r] \& H^1(A_X)\arrow[u,"\varphi_X"']\\
A_Y \& A_X\arrow[l]
\end{tikzcd}
&
\begin{tikzcd}[ampersand replacement=\&,column sep=small,row sep=small]
H^{2m-1}(Y\unsure{.})(m-1)\arrow[r]\arrow[d,"\psi_Y"] \& H^{2n-1}(X)\arrow[d,"\varphi_X"]\\
H^1(B_Y)\arrow[r] \& H^1(B_X)\\
B_Y \& B_X\arrow[l]
\end{tikzcd}\\ \hline
\hfill III & \hfill IV\\
\begin{tikzcd}[ampersand replacement=\&,column sep=small,row sep=small]
H^{2m-1}(Y)(m-1)\arrow[r]\arrow[d,"\psi_Y"] \& H^1(X)\\
H^1(\deleted{\unsure{A}}B_Y)\arrow[r] \& H^1(A_X)\arrow[u,"\varphi_X"']\\
{} \& {}\\[-1ex]
B_Y \& A_X\arrow[l]
\arrow[from=3-1,to=3-2,phantom,"\text{Cas auto-transposé}" description]
\end{tikzcd}
&
\begin{tikzcd}[ampersand replacement=\&,column sep=small,row sep=small]
H^1(Y)\arrow[r,dash,"\unsure{-|\to}" description] \& H^{2n-1}(X)(n-1)\arrow[d,"\psi_X"]\\
H^1(A_Y)\arrow[u,"\varphi_Y"']\arrow[r,dash] \& H^1(B_X)\\
{} \& {}\\[-1ex]
A_Y \& B_X\arrow[l]
\arrow[from=3-1,to=3-2,phantom,"\text{Cas auto-transposé}" description]
\end{tikzcd}\\ \hline
\end{tabular}}

\uline{Cas} I\quad Pour passer des hom de VA à des cycles sur $X\times Y$, il faut inverser $\varphi_Y$

\uline{Cas} II\quad \deleted{\unsure{On peut|Cas pareil}}\hfill --- $\varphi_X$
\margin{$\left.\vphantom{\begin{matrix}1\\1\end{matrix}}\right\}$ il faut des dénominateurs ?}

\uline{Cas} III\quad On peut tjrs passer d'un hom de VA à un cycle : c'est aussi la théorie classique...

\uline{Cas} IV\quad Pour passer des hom de VA à des cycles sur $X\times Y$, il faut inverser $\varphi_Y$ et $\psi_X$, ce qui introduit des dénominateurs.

Pour inverser $\psi_X$, on utilise «th. de Lefschetz». 

Pour inverser $\varphi_X$, on utilise le \unsure{\textquoteleft|\textquotedblleft|\guillemotleft{}}th de Lefschetz'' et le cas particulier des V.A. qui est trivial.
% archive sheet number: 50
\scanpage{52}
1. \uline{Un lemme de théorie des intersections}.

\uline{Th.} 1.1. Soient $X,Y$ projectives \deleted{\unsure{non}} lisses \inserted{connexes,} sur $k$ alg clos, de dim $n$ et $m$

$Z$ un cycle de dimension $m$, codim $n$, sur $X\times Y$.

$A_X,A_Y$ les variétés d'Albanese \unsure{homogènes} de $X,Y$

Alors il existe un homomorphisme de V.A. unique
\[
u_Z:A_Y\longrightarrow A_X
\]
ayant la propriété suivante : pour tout $0$-cycle\margin{de d} $\unsure{\varpi|\mathfrak{a}|\mathfrak{A}}$ sur $Y$, \inserted{de degré $0$,} tel que $Z(\varpi)$ soit défini (c'est \unsure{alors|donc|dans} un $0$-cycle sur $X$), on a
\[
S_X(Z(\varpi))=u(S_Y(\varpi)).
\]
De plus, $u_Z$ ne dépend que de classe d'équiv. \uline{numérique} de $Z$.

\uline{Corollaire} 1.2. Si \deleted{$\unsure{Z}$ \unsure{et} $\varpi$ est \unsure{alb|alg.}-équivalent à zéro,}\inserted{$S_Y(\varpi)=0$} et si $Z(\varpi)$ est défini, \inserted{alors} $S_X(Z(\varpi))=0$ \deleted{\unsure{est} alb-équivalent à $0$.}

\uline{Démonstration}\qquad \uline{Unicité} Elle est à peu près claire, car il existe un ouvert non vide $U$ de $Y$ tel que $y\in U(k)\leadstoZ(y)$ est défini, \deleted{D'autre part,}\inserted{donc $Z(\varpi)$} est défini pour tt $0$-cycle dans $Y$ de support $\subset U$. D'autre part, $\deleted{\unsure{B}}A_Y(k)$ est égal à l'ensemble des $S(\varpi)$, où $\varpi$ est un \deleted{\illegible{1}} $0$-cycle de degré $0$ sur $Y$, de support $\subset U$, donc $u(k):A_Y(k)\unsure{-|\longrightarrow}A_X(k)$ est bien déterminé
% archive sheet number: 51
\scanpage{53}
Pour \uline{l'existence}, considérons l'homomorphisme canonique
\[
\varphi_X:X\longrightarrow A_X^1=X'
\]
et le cycle
\[
Z'=(\varphi_X\times\mathrm{id}_Y)_*(Z)\in\mathrm{Cyc}_m(X'\times Y)
\]
Introduisons également
\[
B=B_X=\underline{\mathrm{Pic}}^{00}_{A_X^1/k}\simeq\underline{\mathrm{Pic}}^0_{A_X/k}
\]
(isomorphe canoniquement à $\underline{\mathrm{Pic}}^{00}_{X/k}$ par l'hom $\underline{\mathrm{Pic}}^0_{A_X^1/k}\unsure{-|\longrightarrow}\underline{\mathrm{Pic}}^0_{X/k}$ transposé de $\varphi_X$). Considérons le cycle $Z''=B\times Z'$ sur $B\times X'\times Y$, et l'élément canonique $\Theta$ de $\mathrm{Pic}(B\times X')$, dont l'image inverse \deleted{\unsure{dans}} $B\times X'\times Y$ \unsure{sera|ser.} notée $\Theta'$. Soit $\deleted{\unsure{\theta|\Theta'}}\theta'$ la classe de diviseurs, mod équivalence numérique, définie par $\Theta'$, et soit de $\overline{\mathrm{m}}$ $z''$ la classe d'équivalence numérique de $Z''$ (qui ne dépend, bien entendu, que de la classe d'équiv. numérique de $Z$). Soit \deleted{$\unsure{\pi:}$}
\[
\pi:B\times X'\times Y\longrightarrow B\times Y
\]
la projection, et formons
\[
c=\pi_*(z''\cdot\unsure{\Theta|\theta})
\]
qui est une classe d'équiv \uline{numérique} de diviseurs sur $B\times Y$, \uline{donc définit} une classe
% circled note
\margin{\unsure{une référence|en référence}}
% archive sheet number: 52
\scanpage{54}
de correspondance divisorielle canonique
\[
c'\in\mathrm{Cl}(B\unsure{,|\times}\deleted{\unsure{A_Y|A}}Y).
\]
\deleted{\unsure{Nous prop.|Nous supposons}} Cette classe définit un homomorphisme \deleted{$Y\longrightarrow\underline{\mathrm{Pic}}^{\unsure{1|0}}_{B/k}$, d'où un homomorphisme} $A_Y\longrightarrow\underline{\mathrm{Pic}}^0_{B/k}$, or par Cartier (bidualité) on a \inserted{un isom can.} $\underline{\mathrm{Pic}}^0_{B/k}\simeq A_X$, donc $c'$ définit un \deleted{\unsure{élément}} homomorphisme de VA
\[
u_Z:A_Y\longrightarrow A_X,
\]
ne dépendant que de la classe d'équivalence numérique $z$ de $Z$. Donc le th. 1 résultera de l'énoncé plus précis :

\uline{Corollaire} 1.3. L'homomorphisme $u_Z:A_Y\unsure{-|\longrightarrow} A_X$ \unsure{qui on|qu'on} vient de construire satisfait à la condition envisagée dans le th. 1.1. \deleted{\unsure{; en|car}}

Cela résultera du

\uline{Lemme} 1.4. Soit $\unsure{\varpi|\mathfrak{a}|\mathfrak{A}}$ un diviseur de degré $0$
\par\noindent\begin{tabular}{@{}|p{.92\linewidth}@{}}
sur $Y$ tel que $ {}^{\deleted{\unsure{t|*}}}Z(\varpi)$ soit défini. Alors l'élément $\Theta(S_X(Z(\varpi)))$ de $\mathrm{Pic}^0(A_X)$ est égal à $c'(\varpi)$.
\end{tabular}
% archive sheet number: 53
\scanpage{55}
\deleted{\unsure{E.|E}} Cela prouvera bien 1.3. car pour prouver $S_X({}^{\deleted{\unsure{t|*}}}Z(\unsure{\varpi|\mathfrak{a}|\mathfrak{A}}))=u_Z(S_Y(\varpi))$ \inserted{\unsure{dans|sur} $A_X(k)$}, il suffit de prouver (grâce à la \deleted{bi}dualité grossière i.e. ensembliste \deleted{$A_X$ \unsure{sur}} déduite de $A_X\xrightarrow{\sim}A_X^{**}=B^*$, \unsure{à savoir} $A_X(k)\longrightarrow\mathrm{Pic}^0(B)$) que $\Theta(S_X({}^{\deleted{\unsure{t|*}}}Z(\varpi)))=\Theta(u_Z(S_Y(\varpi)))$ \inserted{dans $\mathrm{Pic}^0(B)$}. \deleted{\unsure{Or par const.}} Cela résultera de 1.4. et \inserted{de} l'équation
\[
\deleted{\Theta(u_Z(b))=c'(b)\quad\inserted{\text{dans }\deleted{\unsure{B}}\mathrm{Pic}^0(B_X)}\quad\text{pour tout } b\in A_Y(k)}
\]
\[
c'(\varpi)=\Theta(u_Z(S_Y(\varpi)))\qquad\text{dans }\mathrm{Pic}^0(B)
\]
qui résulte de la définition de $u_Z$ en termes de $c'$, cf préliminaires.

D'autre part, 1.4. sera une conséquence immédiate du

\uline{Lemme} 1.5. Soit $\varpi$ comme dans 1.4. Alors il existe un diviseur $D$ dans la classe \unsure{d'|de'}équivalence linéaire de $\Theta$, \unsure{telle|tell|tel} que, \deleted{\unsure{l'on ait}} posant $C=\pi_*(Z''\cdot(D\times Y))\in\deleted{\unsure{Z}}\mathrm{Cyc}^1(B\times Y)$, \deleted{\illegible{1}} on ait
\[
D(\varphi_{X*}({}^{\deleted{\unsure{t|*}}t}Z(\varpi)))=\deleted{\unsure{Z^t|{}^t Z}}C(\varpi)\qquad\text{dans }\mathrm{Cyc}^1(B)
\]
(\inserted{et que} les deux membres de cette égalité soient définis !).
% archive sheet number: 54
\scanpage{56}
La démonstration de 1.5. est de la théorie d'intersections standart, que nous nous dispensons ici d\deleted{\unsure{e}}'expliciter.

Suivant du point de vue cohomologie $l$-adique la construction qu'on vient de faire, on vérifie sans difficulté :

\uline{Théorème} 1.6. Avec les notations du th. 1.1. considérons $\gamma^n(Z)\in H^{2n}(X\times Y)(n)$, et l'homomorphisme correspondant
\[
a(\gamma^n(Z)):H^1(X)\longrightarrow H^1(Y)
\]
\deleted{Considérons d'autre} Considérons d'autre part \deleted{Alors on a} les h\deleted{\unsure{iso|om}}om\inserted{\unsure{o|à}}morphismes canoniques (isom. mod \deleted{\unsure{torsion})}\inserted{groupes finis}) $\varphi_X^*:H^1(A_X)\unsure{-|\longrightarrow} H^1(X)$, $\varphi_Y^*:H^1(A_Y)-H^1(Y)$. On a commutativité dans le diagramme d'homomorphismes
\diagram{\begin{tikzcd}[ampersand replacement=\&]
H^1(X)\arrow[r,"a(\gamma^n(Z))"] \& H^1(Y)\\
H^1(A_X)\arrow[u,"\varphi_X^*"']\arrow[r,"u_Z^{(1)}"] \& H^1(A_Y)\arrow[u,"\varphi_Y^*"']
\end{tikzcd}}
(Les \deleted{\unsure{f}}cohomologies sont les cohomologies à coefficient $\mathrm{Z}_{\ell}$). \deleted{Ce diagramme est \uline{commutatif}}.
% archive sheet number: 55
\scanpage{57}
\deleted{\unsure{Faisant|Passant}} Prenant pour $X,Y$ des VA, on trouve :

\uline{Corollaire} 1.7. \deleted{\unsure{Tout|Soit}} Soient $X,Y$ deux V.A. de dim. $n,m$. Pour qu'un homomorphisme $\psi:H^1(X)\longrightarrow H^1(Y)$ soit définissable par un cycle algébrique $Z$ de \deleted{\unsure{di|dim}} codim $n$ sur $X\times Y$, il faut et il suffit qu'il soit définissable par un homomorphisme de V.A. $u:Y\unsure{-|\longrightarrow}X$ comme $\psi=u^{(1)}$. De plus, $Z$ étant fixé, \deleted{\unsure{ce|u}} $u$ est uniquement déterminé, et indépendant du nb premier $l$ choisi pour les coefficients dans la cohomologie.

\uline{Remarque} 1.8. Généralisation \inserted{de 1.1.} aux schémas projectifs lisses, à fibres \unsure{géom|géom.} connexes, sur une base quelconque $S\deleted{\unsure{Z}}$, en partant d'un $n$-cocycle (au sens du \deleted{\unsure{le}} gradué associé au $K$) $Z\in{\mathrm{CK}}^n(X\times_S Y)$ : on obtient, si $A_Y$ et $A_X$ existent, (plus gén, si $A_Y$ existe
% archive sheet number: 56
\scanpage{58}
et si on se donne un hom de $X$ dans un \deleted{principal homog.} schéma \unsure{paraabélien|parabélien}) sous un schéma abélien $A$) un hom $A_Y\longrightarrow A$, via un élément de ${\mathrm{CK}}^1(A_Y,A^*)$. \inserted{Sur les} fibres \deleted{par fibre}\inserted{\unsure{géom|géom.}}, c'est donné par la construction déjà explicitée\deleted{\unsure{)}}.

2. \uline{Application à l'algébricité de certains homomorphismes de V.A. ou d'\unsure{espaces} de cohomologie $l$-\unsure{adiques|adique}.}

\deleted{Soit $X$ projective lisse connexe sur $k$ alg clos, et soit $z_1\in Z_1^{(\mathrm{num})}(X)$, défini par $Z_1\in\mathrm{Cyc}_1(X)$.}

\uline{Prop.} 2.1. Soient $X,Y$ projectives lisses connexes sur $k$, de dim $n$ resp. $m$, et \deleted{a} $Z_1\in\mathrm{Cyc}_1(X\times Y)$. Alors il existe un homomorphisme
% slanted
\margin{\unsure{echanger|échanger} le rôle de $X$ et $Y$ ?}
unique
\diagram{\begin{tikzcd}[ampersand replacement=\&,column sep=small,row sep=tiny]
u: \& B_X\arrow[r] \& A_Y\\
\& \big\Vert \& \\
\& \underline{\mathrm{Pic}}^{00}_{X/k} \&
\end{tikzcd}}
satisfaisant la condition suivante : pour tout diviseur $D$ sur $X$, algébriquement équiv. à $0$, tel que $ {}^{t}Z(D)=\deleted{\unsure{\pi}}\mathrm{pr}_{2*}((D\times Y)\cdot Z)$ soit défini, on \unsure{a|ait} $u(\mathrm{cl}(D))=S_{\unsure{Y|X}}({}^{t}Z(D))$. Cet homomorphisme
% archive sheet number: 57
